\documentclass[10pt,a4paper]{amsart}
\usepackage[margin=19mm]{geometry}
\usepackage{amsmath,amssymb,mathtools}
\usepackage[scr=rsfs,cal=euler]{mathalfa}
\usepackage{xfrac}
\usepackage[unicode,hidelinks,bookmarks=false]{hyperref}
\newcommand{\ie}{i.e.\ }
\newcommand{\cf}{cf.\ }
\newcommand{\resp}{respectively\ }
\newcommand{\Subsection}{\subsection}
\providecommand{\nobreakdash}{\nobreak\hbox{-}}
\setlength{\emergencystretch}{2em}
\multlinegap0pt
\usepackage{amssymb}
\usepackage{xcolor}
\newtheorem{lemma}{Lemma}
\title{Multiscale Hyperbolic-Parabolic Models for Nonlinear Reactive Transport in Heterogeneously Fractured Porous Media}
\author{Taras Mel'nyk, Sorin Pop, Christian Rohde}
\date{Source: arXiv:2602.16439v1 (CC BY 4.0)}
\begin{document}
\maketitle

\noindent\emph{Source and licence.} Multiscale Hyperbolic-Parabolic Models for Nonlinear Reactive Transport in Heterogeneously Fractured Porous Media. Version arXiv:2602.16439v1 (CC BY 4.0), distributed by arXiv under the Creative Commons Attribution 4.0 International licence, http://creativecommons.org/licenses/by/4.0/, as recorded in the arXiv licence metadata for this exact version and shown on the abstract page https://arxiv.org/abs/2602.16439v1. Full licence notice: \texttt{licenses/arxiv-2602.16439-CC-BY-4.0.txt}.\par\smallskip

\noindent\textit{Editorial scope.} Counted unit: the article's lemma lemma-4-2 (source0.tex lines 1031--1039). The article's complete original proof of it is delivered with it (source0.tex lines 1040--1129, one \texttt{proof} environment of 90 lines). No mathematics is written, repaired or re-derived: the statement and the proof are the article's own lines, in the article's own notation, and no retained line is substituted or re-flowed.  No further source-local context is needed: every cross-reference of every retained line resolves inside the two delivered spans.  Nothing was omitted from any retained span, so the package carries no omission record.  No retained line contains a citation, so the wrapper carries no bibliography and no bibliography entry.  No retained line contains a figure environment, an image-inclusion command or drawing code, so no image is delivered and the package declares no asset dependency.  Editorial formatting only, in addition: an AMS \texttt{amsart} preamble that restates the article's own declarations and macros used by the retained lines, namely the environments \texttt{lemma} on separate counters, together with the standard packages those lines need.  The retained environments print in this document as Lemma 1 in order of appearance, whereas the article prints the same items with the numbering of its own full document.  The complete native source of the article as distributed in the arXiv e-print for this exact version is bundled unchanged as \texttt{sources/194882/source0.tex}.
	\begin{lemma}\label{lemma-4-2}
		Let $\{f_n(t)\}_{n\in \mathbb{N}_0}$ be a sequence of continuous and non-negative functions on $[0,T]$
		such that
		$$
		f_{n+1}(t) \le q\,f_n(t) + C \int_{0}^{t} f_n(\tau)\, d\tau,
		$$
		where $ 0 < q < 1,\ C > 0.$ Then the series $\sum_{n=0}^{\infty} f_n(t) $
		converges uniformly for $t\in[0,T].$ 
	\end{lemma}

	\begin{proof} 
		Let’s define a new majorant sequence:
		\begin{equation}\label{recur-1}
			g_0(t) := M_0 = \max_{t \in [0,T]} f_0(t), \quad
			g_{n+1}(t) = q\,  g_n(t) + C \int_0^t g_n(\tau)\, d\tau.
		\end{equation}
		This sequence satisfies the same recurrence as $\{f_n\},$ and it is easy to verify that
		$f_n(t) \le g_n(t)$ for $t\in [0,T].$
		
		Now let us find a solution of  the linear recurrence \eqref{recur-1}. To do this we define the bivariate generating function
		\[
		G(x,t)\;=\;\sum_{n=0}^{\infty}g_{n}(t)\,x^{n}.
		\]
		At \(t=0\), the integral vanishes and \(g_{n}(0)=q\,g_{n-1}(0)\) with \(g_{0}(0)=M_{0}\).  By induction
		$g_{n}(0)=M_{0}\,q^{n},$
		so
		\[
		G(x,0)
		=\sum_{n=0}^{\infty}M_{0}\,q^{n}\,x^{n}
		=\frac{M_{0}}{1-qx} \quad \text{for} \ \ x \in [0, \tfrac{1}{q}).
		\]
		
		Multiplying recurrence \eqref{recur-1} by $x^{n}$ and summing over $n\ge0,$ we get
		functional equation for $G$:
		\begin{equation}\label{recur-2}
			\frac{G(x,t)-M_{0}}{x}
			=q\,G(x,t)
			+ C\int_{0}^{t}G(x,\tau)\,d\tau
			\;\Longrightarrow\;
			(1-qx)\,G(x,t)
			= M_{0} + C\,x\int_{0}^{t}G(x,\tau)\,d\tau.
		\end{equation}
		Differentiating \eqref{recur-2} with respect to $t,$ we obtain 
		\[
		\frac{\partial G}{\partial t}
		= \frac{C\,x}{1-qx}\,G,
		\]
		whence 
		\[
		G(x,t)
		=G(x,0)\,
		\exp\Bigl(\frac{C\,x\,t}{1-qx}\Bigr)
		=\frac{M_{0}}{1-qx}
		\exp\Bigl(\frac{C\,x\,t}{1-qx}\Bigr).
		\]
		
		Using obvious expansions
		\[
		\frac{1}{1-qx}
		=\sum_{m=0}^{\infty}q^{m}x^{m},
		\quad
		\exp\Bigl(\frac{C\,x\,t}{1-qx}\Bigr)
		=\sum_{k=0}^{\infty}\frac{1}{k!}
		\Bigl(\frac{C\,t\,x}{1-qx}\Bigr)^{k}
		=\sum_{k=0}^{\infty}\frac{(C\,t)^{k}}{k!}\,x^{k}
		\sum_{j=0}^{\infty}\binom{k+j-1}{j}(q\,x)^{j},
		\]
		we deduce
		\[
		G(x,t)
		=M_{0}
		\sum_{m,k,j\ge0}
		q^{m}\,
		\frac{(C\,t)^{k}}{k!}\,
		\binom{k+j-1}{j}\,
		q^{\,j}\,
		x^{\,m+k+j}.
		\]
		The coefficient of $x^{n}$ is
		\[
		g_{n}(t)
		= M_{0}
		\sum_{\substack{m,k,j\ge0\\m+k+j=n}}
		q^{\,m+j}
		\frac{(C\,t)^{k}}{k!}
		\binom{k+j-1}{j}.
		\]
		Collapsing  the triple sum, we have
		\[
		g_{n}(t)
		= M_{0}
		\sum_{k=0}^{n}
		\binom{n}{k}
		q^{\,n-k}
		\frac{(C\,t)^{k}}{k!}.
		\]
		
		Taking into account that $q \in (0, 1),$ we verify with the  Weierstrass $M$-test that the series $\sum_{n=0}^{\infty} g_n(t) $
		converges uniformly for $t\in[0,T].$
	\end{proof}

\end{document}
